% =============================================
% C4P 数理モデル（マスター版）
% 書籍『クリエイト・フォー・パーソナリティ』収録の数理モデル章（ver3）を出発点とし、
% 以後の改訂をこのファイルで行う。
% 書籍版からの差分は各節冒頭の【書籍版との差分】注記で明示する。
% エンジン: lualatex（フォント指定なし・素のjlreqでコンパイル可能）
% =============================================

\documentclass[a4paper, fontsize=10pt]{jlreq}

\usepackage{amssymb}
\usepackage{amsmath}
\usepackage{bm,mathtools}
\usepackage{booktabs}
\usepackage{enumitem}
\usepackage{tikz}
\usetikzlibrary{arrows.meta, positioning, decorations.pathreplacing, calc}

% ---- Notation（書籍版から継承） ----
\newcommand{\R}{\mathbb{R}}
\newcommand{\vx}{\bm{x}}
\newcommand{\va}{\bm{a}}
\newcommand{\vs}{\bm{s}}
\newcommand{\vu}{\bm{u}}
\newcommand{\norm}[1]{\left\lVert #1 \right\rVert}
\newcommand{\A}{\mathcal{A}}
\newcommand{\Si}{\mathcal{S}_i}
\newcommand{\Sj}{\mathcal{S}_j}
\newcommand{\Vi}{\mathcal{V}_i}
\newcommand{\Vj}{\mathcal{V}_j}
\newcommand{\Phii}{\Phi_i}
\newcommand{\Phij}{\Phi_j}
\newcommand{\Psii}{\Psi_i}
\newcommand{\Psij}{\Psi_j}
\newcommand{\Ci}{\mathcal{C}_i}
\newcommand{\Fi}{\mathcal{F}_i}
\newcommand{\I}{\mathcal{I}}
\newcommand{\U}{\mathcal{U}}

% ---- 本稿の新規記号 ----
\newcommand{\dV}{\bm{\delta}^{V}}          % 選好的意図（V_i上）
\newcommand{\dS}{\bm{\delta}^{S}}          % 記述的意図（S_i上）
\newcommand{\Di}{\mathcal{D}_i}            % 記述候補集合
\newcommand{\Ri}{\mathcal{R}_i}            % 到達可能集合
\newcommand{\traj}{\omega}                 % 軌跡

% ---- Q&A用カウンタ ----
\newcounter{qacounter}
\newcommand{\qa}[1]{%
  \stepcounter{qacounter}%
  \vspace{1.5\baselineskip}%
  \noindent{\bfseries Q\theqacounter. #1}\par
  \vspace{0.5\baselineskip}%
}

% ---- 仮定用カウンタ ----
\newcounter{assumption}
\newcommand{\assumption}[1]{%
  \stepcounter{assumption}%
  \par\medskip\noindent\textbf{仮定A\theassumption}\quad #1\par\medskip
}

% ---- 差分注記（書籍版との関係を各節冒頭に明示） ----
\newcommand{\diffstatus}[1]{%
  \par\noindent{\small【書籍版との差分】#1}\par\medskip}

\widowpenalty=10000
\clubpenalty=10000

\title{C4P数理モデル}
\author{有村崚}
\date{2026-07-03}

\begin{document}
\maketitle

% ============================================================
\section*{書籍収録版（ver3）からの変更点}
% ============================================================

本稿は書籍『クリエイト・フォー・パーソナリティ』収録の
数理モデル章（以下、書籍版）を出発点とし、以後の改訂を加えていくマスター版である。
書籍版の本文はすべて本稿に含まれており、
各節の冒頭に【書籍版との差分】として変更の有無と内容を明示する。

変更は「書籍版が自分で認めていた曖昧さの解消」と
「本文にあってモデルに無かった概念の取り込み」の二系統からなる。

\paragraph{系統1：書籍版の自己申告課題への対応}
\begin{enumerate}[leftmargin=2em]
  \item \textbf{意図の二層化}（第\ref{subsec:intent}節）——
    書籍版「本モデルの射程」節で「意図の所在は自明ではない」と保留していた問題に決着をつける。
    意図を選好的意図$\dV \in \Vi$と記述的意図$\dS \in \Si$に分離し、
    創作を二段階の逆問題として書き直す。
    本文「ベースの音ってどれですか」の認知／価値判断の二フェーズ分離と正確に対応する。
  \item \textbf{仮定の明示化}（第\ref{sec:assumptions}節）——
    書籍版では$\Vi$のノルム仮定や$\Pi_i$の1ステップ不変性が本文中に散在していた。
    冒頭に仮定リストとして集約し、「都合の悪い部分が浮き上がる」効能を強化する。
  \item \textbf{状態量と経路量の区別}（第\ref{sec:state-path}節）——
    C4Q/C4Pの分離を「マキシマイズか否か」ではなく
    「状態の関数か、軌跡の汎関数か」という数学的な型の違いとして再定義する。
    これにより、異時点の$\bm{\delta}$を同一空間で比較できないという書籍版の残課題は、
    「比較を要求しない定義」に置き換わることで解消される。
  \item \textbf{選択カーネル$\sigma_i$の昇格}（第\ref{subsec:kernel}節）——
    書籍版では「射程」節の脚注的言及に留まっていた選択確率分布を本文に昇格し、
    パーソナリティの構成要素として正式に位置づける。
\end{enumerate}

\paragraph{系統2：本文にあってモデルに無かった概念の定式化}
\begin{enumerate}[leftmargin=2em]\setcounter{enumi}{4}
  \item \textbf{完成の定式化}（第\ref{sec:completion}節）——
    本文「完成宣言」。完成を収束ではなく停止の選択として定式化し、
    停止パターンをパーソナリティの第4の構成要素に加える。
  \item \textbf{凡庸さの定式化}（第\ref{sec:mediocrity}節）——
    本文「脱凡庸」。同質化を到達可能集合$\Ri$の個人間重なりとして記述し、
    「技術向上はパーソナリティの出力経路を増やす」という本文の主張に構造を与える。
  \item \textbf{観測の交換}（第\ref{sec:exchange}節）——
    本文「馴れ合いしようや」。他者の反応の二つの使い途
    （$Y$の代理指標として／自分の$\Phii, \Psii$の更新入力として）を分離し、
    馴れ合いの3分類をモデル上で判別可能にする。
  \item \textbf{運の所在}（第\ref{sec:luck}節）——
    本文「不運な人」「アウトオブスコープ」。モデルの外生変数を列挙し、
    「パーソナリティだけが運に対して平等」という主張の正確な意味
    （定義可能性の環境非依存性）を与える。
\end{enumerate}

Q\&Aは書籍版の30問（Q1--Q30）をすべて収録し、
新設概念に対応する6問（Q31--Q36）を追加した。
引き継いだ30問は、記号と参照先を本稿に合わせて更新した以外は書籍版のままである。

\clearpage

% ============================================================
\section*{記法ガイド}
% ============================================================

本稿では数学の記号を使って議論を進める。
目的は数学をすることではなく、態度を正確に記述することである。
以下の記号の読み方がわかれば本稿は追える。

\bigskip

{\small
\begin{tabular}{@{}llp{70mm}@{}}
\toprule
記号 & 読み & 本稿での意味 \\
\midrule
$\A$ & エー & 属性空間。創作物の物理的状態すべてが張る空間 \\
$\Si$ & エスアイ & 意味空間。個人$i$の知覚・認知が畳み込んだ空間 \\
$\Vi$ & ブイアイ & 選好潜在空間。好き嫌いの座標系 \\
$\Phii$ & ファイ & 知覚・認知写像。何を知覚できるか \\
$\Psii$ & プサイ & 評価・選好写像。何を好むか \\
$\Pi_i$ & パイ & 合成射影（$= \Psii \circ \Phii$）。パーソナリティの発現部 \\
$\dV$ & デルタブイ & \textbf{選好的意図}。「もっと好きな感じにしたい」（$\Vi$上）\\
$\dS$ & デルタエス & \textbf{記述的意図}。「音色をこう変えたい」（$\Si$上）\\
$\Di$ & ディーアイ & \textbf{記述候補集合}。$\dV$を実現する$\dS$の選択肢 \\
$\Ci$ & シーアイ & 候補集合。$\dS$を実現する操作の選択肢 \\
$\Fi$ & エフアイ & 実行可能領域。技術的に実行できる操作の範囲 \\
$\sigma_i$ & シグマ & \textbf{選択カーネル}。候補からどれを選ぶかの傾向 \\
$\Ri(\va)$ & アール & \textbf{到達可能集合}。$\va$から1ステップで行ける状態の集合 \\
$\tau$ & タウ & \textbf{停止時刻}。「完成」と宣言する時点 \\
$r_{j \to i}$ & アール & \textbf{観測の交換}。$j$の反応の$i$への報告 \\
$U_\Phi, U_\Psi$ & ユー & \textbf{写像の更新規則}。変容の抽象的な記述 \\
$\U$ & ユー（飾り文字） & 操作空間。$\vu$が属する空間 \\
$T$ & ティー & 状態遷移写像。操作を状態に適用する \\
$Y(\va)$ & ワイ & 社会的クオリティ。他者の反応の集計量 \\
$\va$ & エー（太字） & 創作物の状態（$\A$の元） \\
$\vu$ & ユー（太字） & 操作（候補集合から選ぶもの） \\
$\traj$ & オメガ & \textbf{軌跡}。$(\va, \dV, \dS, \vu)$の系列全体 \\
$\bm{\epsilon}_i$ & イプシロン & 実現誤差。意図と結果のずれ（本文中では$\lVert\epsilon\rVert$と略記） \\
\bottomrule
\end{tabular}
}

\bigskip

太字は本稿で新設または昇格した記号である。
書籍版の意図ベクトル$\bm{\delta}_i$は、本稿では選好的意図$\dV$に対応する。

$\Psii \circ \Phii$ は「まず $\Phii$（知覚フィルタ）を通し、次に $\Psii$（選好フィルタ）を通す」という意味である。
二つの写像を順番に適用しているだけである。

\clearpage

書籍『クリエイト・フォー・パーソナリティ』で提唱したC4Q／C4Pの概念に数理的な定式化を与える。
本文の「ベースの音ってどれですか」章で述べた「認知フェーズ」「価値判断フェーズ」は、
本稿における知覚写像$\Phii$と選好写像$\Psii$に対応する。

本稿では、C4Qを\textbf{2つのレイヤー}に分解する。
\begin{itemize}[leftmargin=2em]
  \item \textbf{レイヤー1：技術}（個人内で完結）——
    自分の意図を属性空間$\A$上で正確に実現できる度合い。
    実行可能領域$\Fi$の広さと実現精度に対応する。
  \item \textbf{レイヤー2：社会的クオリティ}（他者の集合にわたる集計）——
    他者の反応の総和$Y(\va)$。
    C4Qが追求する「客観的な凄さ」とはこれのことである。
\end{itemize}

% ============================================================
\section{仮定}\label{sec:assumptions}
% ============================================================

\diffstatus{新設。書籍版では本文中に散在していた作業仮説をここに集約した。}

以降のすべての式はこの仮定リストの上に立っており、
仮定が崩れる場所ではモデルも黙る。

\assumption{$\A$は十分高次元の実ベクトル空間または関数空間の部分集合である。}

\assumption{各時点$t$において、$\Si^{(t)}$と$\Vi^{(t)}$は実ベクトル空間であり、
$\Vi^{(t)}$にはノルム$\norm{\cdot}$が入る。
この構造は\textbf{同一時点の近傍でのみ}使用する。
異なる時点の$\Vi^{(t)}$と$\Vi^{(t')}$の間の同一視は仮定しない。}

\assumption{創作の1ステップの時間幅は、$\Phii$と$\Psii$の変容の時定数より十分短い。
すなわち1ステップの内部では$\Pi_i$を固定とみなす。
変容（第\ref{sec:dynamics}節）は多数ステップにわたる長い時間スケールで扱う。}

\assumption{$T : \A \times \U \to \A$は与えられた道具立て（楽器・DAW・絵筆・AI）で決まる写像であり、
個人には属さない。個人に属するのは$\Fi \subset \U$と選択である。}

\assumption{$\Phii, \Psii$の具体的な関数形は与えない。
本モデルの目的は数値化ではなく態度の正確な記述である。}

仮定A2が本稿の要である。
書籍版は「異なる時点の$\bm{\delta}_i$を同じ空間上で比較できない」ことを
未解決の課題として残していたが、
本稿ではノルムの使用を同一時点の近傍に制限し、
時点をまたぐ量はすべて「比較」ではなく「蓄積」として定義する
（第\ref{sec:state-path}節）。
比較できないものを比較しない形に定義を直すのが、本稿の基本方針である。

% ============================================================
\section{三層射影モデル}\label{sec:three-layers}
% ============================================================

\diffstatus{変更なし。書籍版の第1節をそのまま収録する。}

創作物が個人によって鑑賞・評価されるプロセスを、
以下の三つの空間の間の射影として記述する（図\ref{fig:three-layers}）。

\begin{enumerate}
  \item \textbf{属性空間} $\A$：
    創作物が物理的に取りうるあらゆる状態が張る、バカでかい空間。
  \item \textbf{意味空間} $\Si$：
    個人$i$の知覚・認知・知識によって$\A$から畳み込まれる、ややコンパクトな空間。
  \item \textbf{選好潜在空間} $\Vi$：
    意味空間$\Si$からさらに好悪の軸へ射影された空間。
\end{enumerate}

射影の連鎖は
\[
  \va \in \A
  \xrightarrow{\;\Phii\;}
  \vs_i \in \Si
  \xrightarrow{\;\Psii\;}
  \vx_i \in \Vi
\]
であり、合成射影$\Pi_i = \Psii \circ \Phii$として表される。

\begin{figure}[htbp]
\centering
\begin{tikzpicture}[
  >=Stealth,
  box/.style={draw, rounded corners, minimum width=48mm, minimum height=12mm,
              align=center, font=\footnotesize},
  arr/.style={->, thick, font=\footnotesize},
]
  \node[box] (A) {属性空間 $\A$\\（物理状態・実質無限次元）};
  \node[box, below=14mm of A] (S) {意味空間 $\Si$\\（知覚・認知・知識による畳み込み）};
  \node[box, below=14mm of S] (V) {選好潜在空間 $\Vi$\\（好悪・お気持ちの座標系）};

  \draw[arr] (A) -- node[right, xshift=2mm]{$\Phii$：知覚・認知写像} (S);
  \draw[arr] (S) -- node[right, xshift=2mm]{$\Psii$：評価・選好写像} (V);

  \draw[decorate, decoration={brace, amplitude=6pt, mirror}, thick]
    ([xshift=38mm]A.north) -- ([xshift=38mm]V.south)
    node[midway, right=8pt, align=left, font=\footnotesize]
    {$\Pi_i = \Psii \circ \Phii$\\"パーソナリティ"\\の発現部};
\end{tikzpicture}
\caption{三層射影モデルの概観。$\Pi_i$が二つの写像に分解される。}
\label{fig:three-layers}
\end{figure}

% ------------------------------------------------------------
\subsection{属性空間 $\A$}
% ------------------------------------------------------------

創作物の物理的状態の全体が張る空間を\textbf{属性空間}と呼び、$\A$と書く。
次元は実質的に無限であり、$\A \subset \R^{D}$（$D \to \infty$）
または適切な関数空間として定式化する。

デジタル形式の音楽であれば、波形のサンプル値の列がこの空間の一点を定める。
例えば、CDに収録された3分間のステレオ音源（44.1\,kHz）は
$D = 44100 \times 2 \times 180 = 15{,}876{,}000$次元のベクトルとして表現できる。
各次元の値域はビット深度で決まり（16ビットなら$0$--$65535$）、次元数には影響しない。
画像であればピクセル値の列、
プロダクトであれば寸法・材質・色彩・重量など
計測可能なあらゆる物性値が成分となる。

属性空間は創作物の側に属する空間であり、人間に依存しない。
同じ楽曲は誰が聴いても（物理的には）同じ波形を耳に届ける。

ただし$\A$は巨大だが、実際に存在する（あるいは実現可能な）創作物は
$\A$の中の極めて薄い部分多様体上に乗っている。
ランダムなサンプル値の列はホワイトノイズであり、
「楽曲」と認識される波形は$\A$のごく限られた領域に集中する。

% ------------------------------------------------------------
\subsection{意味空間 $\Si$}
% ------------------------------------------------------------

個人$i$の知覚・認知系を写像
\begin{equation}
  \Phii : \A \longrightarrow \Si
\end{equation}
として表す。$\Si$は個人$i$にとって知覚・認識可能な意味的特徴が張る空間であり、
\textbf{意味空間}と呼ぶ。

$\Phii$は大規模な次元削減である。
44.1\,kHzの波形データが、
「テンポ130」「マイナーキー」「リバーブが深い」
「ラテンのリズムが用いられている」
といった意味的特徴の組に変換される。
この変換は個人に依存する：

\paragraph{例1（和声の訓練）}
和声の訓練を受けた聴取者$i$の$\Si$には
「トライトーンの解決」「裏コードの使用」「ドミナントモーション」
といった次元が存在する。
訓練を受けていない聴取者$j$の$\Sj$には
これらの次元がそもそも存在しないか、
あったとしても極めて粗い解像度でしか符号化されない。
同じ波形$\va \in \A$に対して、$\Phii(\va)$と$\Phij(\va)$は
異なる次元数の異なるベクトルになりうる。

\paragraph{例2（文化的背景）}
日本の演歌を聴いたとき、日本語話者は歌詞の意味内容を
$\Si$の複数の次元として符号化するが、
日本語を解さない聴取者はその次元を持たず、
代わりに音韻的・音響的特徴のみを意味空間に持つ。
しかし後者が旋律やこぶしの構造に対して持つ解像度が
前者より高い場合もありえる。$\Si$の「良し悪し」は一意には定まらない。

\paragraph{$\Phii$の性質}
$\Phii$は一般に以下の性質を持つ：
\begin{itemize}[leftmargin=2em]
  \item \textbf{非単射}：
    $\A$上で異なる二つの状態$\va, \va'$が$\Phii(\va) = \Phii(\va')$
    となりうる（人間が区別できない物理的差異は意味空間で潰れる）
  \item \textbf{非線形}：
    知覚は一般に非線形である
  \item \textbf{時変}：
    学習や経験により$\Phii$は変化する。
    新しい概念を獲得すると$\Si$の次元が増え、複雑になる
  \item \textbf{文脈依存}：
    同じ$\va$でも、先行する聴取経験や状況によって
    $\Phii(\va)$の値が変わりうる
\end{itemize}

% ------------------------------------------------------------
\subsection{選好潜在空間 $\Vi$}
% ------------------------------------------------------------

個人$i$の選好構造を写像
\begin{equation}
  \Psii : \Si \longrightarrow \Vi
\end{equation}
として表す。$\Vi$は個人$i$が意味空間上の特徴に対して付与する
好悪・魅力・嫌悪などの評価が張る空間であり、
\textbf{選好潜在空間}と呼ぶ。

すなわち
\begin{equation}
  \vx_i \;=\; \Psii\bigl(\Phii(\va)\bigr)
      \;=\; (\Psii \circ \Phii)(\va)
\end{equation}
である。

\paragraph{例（意味空間と選好の分離）}
とあるメロディを聴いた時、「これはメジャースケールである」という認知は$\Si$のレベルで起きるが、
それが好きか嫌いかは$\Vi$のレベルの話である。
メジャースケールを認知できる二人の聴取者が、
一方はそれに強い魅力を感じ（$\Vi$で高い値）、
他方は特に何も感じない（$\Vi$でゼロ付近）ということがありえる。
$\Si$が同じでも$\Vi$は異なる。

% ------------------------------------------------------------
\subsection{合成射影とパーソナリティの所在}
% ------------------------------------------------------------

三層モデルでは$\Pi_i = \Psii \circ \Phii$として分解されるため、
パーソナリティ（個人性）は$\Phii$と$\Psii$の\textbf{両方}に宿る。

\begin{itemize}[leftmargin=2em]
  \item $\Phii$に宿るパーソナリティ：
    「何を知覚・認知できるか」「どのような概念的語彙を持つか」
    ——\textbf{知覚的パーソナリティ}
  \item $\Psii$に宿るパーソナリティ：
    「認知した特徴のうち何を好み、何を嫌うか」
    ——\textbf{選好的パーソナリティ}
\end{itemize}

よさの定義権が個々人に宿るということは、
$\Phii$（何を観点として認知するか）と
$\Psii$（認知した特徴をどう評価するか）の
両方が個人に属する、ということの言い換えである。

% ============================================================
\section{創作行為の定式化：候補集合モデル}\label{sec:creative-act}
% ============================================================

\diffstatus{改訂。意図を$\dV$と$\dS$に二層化し（第\ref{subsec:intent}節）、
候補集合を二段階に分解した（第\ref{subsec:candidate-set}節）。
創作の1ステップは4選択肢から5ステップ構成に改まり、選択点が2箇所になった。
選択カーネル$\sigma_i$を新設した（第\ref{subsec:kernel}節）。
「鑑賞と創作の方向」「実行可能領域と技術」「意図の非定常性」は書籍版のままである。}

% ------------------------------------------------------------
\subsection{鑑賞と創作の方向}
% ------------------------------------------------------------

鑑賞は$\Pi_i$の順方向適用として定式化される。
本節では創作を逆方向の操作として位置づける。

\begin{itemize}[leftmargin=2em]
  \item \textbf{鑑賞}：所与の$\va \in \A$に対して、
    $\Pi_i(\va) = \vx_i \in \Vi$を構成する（順方向）
  \item \textbf{創作}：$\Vi$上の意図に基づいて、
    $\A$上の状態を変化させる（逆方向）
\end{itemize}

鑑賞が「$\A$の住人を$\Vi$の言葉に翻訳する行為」であるならば、
創作は「$\Vi$の言葉を$\A$の住人に翻訳し返す行為」である。

ただし、$\Pi_i$は非単射であるため、
\textbf{厳密な逆写像$\Pi_i^{-1}$は存在しない}。
$\Vi$上の同じ一点に対して、$\A$上の状態は一般に複数対応する。
$\A$や$\U$の自由度は$\Vi$に比べて圧倒的に大きい
（1500万次元の波形空間に対して好悪の座標系は遥かに低次元である）ため、
$\Pi_i \circ T$は大規模な次元圧縮であり、前像は必然的に多対一になる。
この非一意性が創作行為の構造の核心であり、
創作にパーソナリティが発現する数理的な根拠となる。

% ------------------------------------------------------------
\subsection{意図の二層構造}\label{subsec:intent}
% ------------------------------------------------------------

創作者$i$が時刻$t$における制作途中の状態$\va^{(t)} \in \A$を持つとする。
創作者は自作を鑑賞し（$\Pi_i$の順方向適用）、
\begin{equation}
  \vx_i^{(t)} = \Pi_i\bigl(\va^{(t)}\bigr) \in \Vi
\end{equation}
を得る。これは「今の制作物を自分のフィルタで聴いた/見た結果」である。

このとき創作者が持つ「変化への意図」を、書籍版は単一のベクトル$\bm{\delta}_i \in \Vi$と置いたが、
「射程」節で自ら認めていたとおり、
「もう少し暗くしたい」が$\Vi$上の判断なのか$\Si$上の操作指示なのかは曖昧だった。
本稿ではこれを曖昧さではなく\textbf{構造}として書く。
意図は二層ある。

\begin{itemize}[leftmargin=2em]
  \item \textbf{選好的意図} $\dV_i \in \Vi$：
    「もっと好きな感じにしたい」「ここが足りない」「ここは違う」。
    好悪の座標系上の変位。方向（「暗くしたい」）と大きさ（「どれくらい」）を持つが、
    作品を\textbf{どう}変えるかの情報は含まない。
  \item \textbf{記述的意図} $\dS_i \in \Si$：
    「音色を丸くしたい」「テンポを落としたい」。
    意味空間上の変位。作品をどう変えるかを、
    本人の語彙（$\Si$の次元）で記述したもの。
\end{itemize}

創作の駆動力は$\dV$であり、$\dS$は$\dV$を実現するための中間表現である。
「暗くしたい（$\dV$）」に対して
「マイナーコードにする」「テンポを落とす」「高域を削る」（それぞれ$\dS$）が候補になる、
という関係である。

% ------------------------------------------------------------
\subsection{二段階の候補集合}\label{subsec:candidate-set}
% ------------------------------------------------------------

意図の二層化により、創作は二段階の逆問題として書ける。

\paragraph{第1段（$\Vi \to \Si$）：記述候補集合}
選好的意図$\dV$を実現する記述的意図の集合を
\begin{equation}\label{eq:desc-set}
  \Di(\dV,\, \vs)
  \;=\;
  \bigl\{\, \dS \in \Si
    \;\bigm|\;
    \Psii(\vs + \dS) - \Psii(\vs) \approx \dV
  \,\bigr\}
\end{equation}
と定義する。「好きな感じに近づく意味的変更の選択肢」である。

\paragraph{第2段（$\Si \to \A$）：操作候補集合}
記述的意図$\dS$を実現する操作の集合を
\begin{equation}\label{eq:op-set}
  \Ci(\dS,\, \va)
  \;=\;
  \bigl\{\, \vu \in \U
    \;\bigm|\;
    \Phii(T(\va,\vu)) - \Phii(\va) \approx \dS
  \,\bigr\}
\end{equation}
と定義する。「その意味的変更を物理的に起こす操作の選択肢」である。
ここで$\U$は\textbf{操作空間}であり、
$\vu \in \U$は「波形を加工する」「テキストを編集する」「粘土を削る」
といった創作上の操作を表す。
$T : \A \times \U \to \A$は\textbf{状態遷移写像}であり、
状態$\va$に操作$\vu$を施した結果の新しい状態を返す。
$\U$と$\A$は一般に異なる空間である。

なお、式\eqref{eq:desc-set}\eqref{eq:op-set}の$\approx$は意図的に曖昧にしてある。
仮定A2のノルムを使えば角度条件や円錐近傍で形式化できるが、
具体的な閾値条件を課すのは時期尚早である。
「だいたいその方向にだいたいその分だけ動く選択肢の集まり」
という直感で十分である。
$T$や$\U$の厳密な代数的構造も同様に、本稿では追求しない。

\paragraph{例（音楽制作）}
「もう少し暗い雰囲気にしたい」
（$\dV$が$\Vi$上で「暗さ」の方向を指している）
に対して、
\begin{itemize}[leftmargin=2em]
  \item マイナーコードに差し替える
  \item テンポを落とす
  \item リバーブを深くする
  \item 高域をカットする
  \item 歌詞の語彙を変える
\end{itemize}
これらはすべて$\Di(\dV,\, \vs)$の要素である。
意味空間上では全く異なる変更だが、
$\Psii$を通して見たときに「暗さの方向に動く」という点で共通する。
さらに、そのうちの一つ「リバーブを深くする」（$\dS$）を選んだとしても、
どのリバーブをどのパラメータで、どのトラックにかけるかという
$\Ci(\dS,\, \va)$からの選択が残っている。

\paragraph{書籍版との関係}
書籍版の候補集合
$\Ci(\bm{\delta}, \va) = \{\vu \mid \Pi_i(T(\va,\vu)) - \Pi_i(\va) \approx \bm{\delta}\}$
は、二段階を合成して1本にしたものに相当する。
分解して得られるものは次の観察である：
\textbf{パーソナリティの発現点が2箇所に増える。}
同じ$\dV$（暗くしたい）に対して
$\Di$からどの$\dS$を選ぶか（和声で行くか、音色で行くか）が第1の選択、
同じ$\dS$（音色を丸くする）に対して
$\Ci$からどの$\vu$を選ぶか（EQを削るか、マイクを替えるか）が第2の選択である。
本文「ベースの音ってどれですか」の認知／価値判断の二フェーズが、
鑑賞（順方向）だけでなく創作（逆方向）にも対称的に現れる。

% ------------------------------------------------------------
\subsection{候補集合からの選択としての創作}\label{subsec:creative-step}
% ------------------------------------------------------------

$\Di$も$\Ci$も一般に複数の要素を持つ。
これが創作行為の本質的な構造である。

創作の1ステップを以下で定義する（図\ref{fig:loop}）：
\begin{enumerate}[leftmargin=2em]
  \item \textbf{自作の鑑賞}：$\vx_i^{(t)} = \Pi_i(\va^{(t)})$を得る
  \item \textbf{選好的意図の認知}：$\dV{}^{(t)} \in \Vi$を得る
  \item \textbf{記述の選択}：$\Di(\dV{}^{(t)},\, \vs^{(t)})$から$\dS{}^{(t)}$を選ぶ
  \item \textbf{操作の選択}：$\Ci(\dS{}^{(t)},\, \va^{(t)}) \cap \Fi$から$\vu^{(t)}$を選ぶ
  \item \textbf{更新}：$\va^{(t+1)} = T(\va^{(t)},\, \vu^{(t)})$
\end{enumerate}

\textbf{ステップ3と4の選択は候補集合の構造だけからは決まらない。
どれを選ぶかが、創作におけるパーソナリティの発現である。}

候補集合の構成が「暗黙的」であることに注意する。
実際の創作者が候補を意識的に列挙しているわけではない。
ギターを歪ませたい時、「なんとなくいい感じになるようにゲインのつまみを触る」のであって、
ギタリストは「候補を並べて比較して選ぶ」わけではない。
モデルが記述しているのは行為の構造であって体験ではない。
リンゴがニュートンの微分方程式を解いているわけではないのと同じである。

\begin{figure}[htbp]
\centering
\begin{tikzpicture}[
  >=Stealth,
  box/.style={draw, rounded corners, minimum width=46mm, minimum height=12mm,
              align=center, font=\footnotesize},
  arr/.style={->, thick, font=\scriptsize},
]
  \node[box] (a) {状態 $\va^{(t)}$};
  \node[box, right=16mm of a] (x) {鑑賞 $\vx^{(t)} = \Pi_i(\va^{(t)})$};
  \node[box, below=10mm of x] (dv) {選好的意図 $\dV{}^{(t)}$};
  \node[box, below=10mm of dv] (ds) {記述的意図 $\dS{}^{(t)}$\\（$\Di$から選択）};
  \node[box] (anew) at (a.center |- dv.center) {状態 $\va^{(t+1)}$\\（$= T(\va^{(t)}, \vu^{(t)})$）};
  \node[box] (u) at (a.center |- ds.center) {操作 $\vu^{(t)}$\\（$\Ci \cap \Fi$から選択）};

  \draw[arr] (a) -- node[above]{鑑賞（順方向）} (x);
  \draw[arr] (x) -- (dv);
  \draw[arr] (dv) -- node[right]{第1の選択} (ds);
  \draw[arr] (ds) -- node[above]{第2の選択} (u);
  \draw[arr] (u) -- node[left]{$T$の適用} (anew);
  \draw[arr, dashed] (anew) -- node[left]{次ステップへ} (a);
\end{tikzpicture}
\caption{創作ループ。選択点が2箇所ある。}
\label{fig:loop}
\end{figure}

% ------------------------------------------------------------
\subsection{選択カーネル}\label{subsec:kernel}
% ------------------------------------------------------------

書籍版で脚注的に触れていた選択確率分布を本文に昇格する。
ステップ3・4の選択の傾向を、条件付き分布
\begin{equation}
  \dS \sim \sigma_i^{D}\bigl(\cdot \mid \dV,\, \vs\bigr),
  \qquad
  \vu \sim \sigma_i^{C}\bigl(\cdot \mid \dS,\, \va\bigr)
\end{equation}
として表し、まとめて\textbf{選択カーネル}$\sigma_i$と呼ぶ。
台（サポート）はそれぞれ$\Di$と$\Ci \cap \Fi$に含まれる。

$\sigma_i$は「同じ意図に対していつもどう動くか」の癖であり、
パーソナリティの構成要素の一つである（第\ref{subsec:personality-redef}節）。
分布として書くのは確率的行動を主張するためではなく、
決定則（デルタ分布）も含む最も弱い形で「傾向」を書けるからである。
関数形は与えない（仮定A5）。

$\sigma_i$の導入で、本文「まねっこ」やパクリの議論が1行で書ける：
\textbf{複製できるのは$\va$であって$\sigma_i$ではない。}
完全なコピー品が存在しても、
それを生成した選択の傾向はコピー元にしか属さない。

% ------------------------------------------------------------
\subsection{実行可能領域と技術}\label{subsec:feasible}
% ------------------------------------------------------------

現実には、候補集合のすべての要素が実行可能ではない。
個人$i$の技術的能力に対応する\textbf{実行可能領域}$\Fi \subset \U$を導入すると、
実際に選択可能な操作は
\begin{equation}
  \Ci(\dS,\, \va) \;\cap\; \Fi
\end{equation}
に制限される。

例えば「ギターのフレーズが頭の中では鳴っているが指が追いつかない」とか
「この色をこう塗りたいが技法を知らない」みたいなのは、
「やりたいことは分かっているが、$\U$上の操作ができない」状態と言える。
$\Ci(\dS,\, \va) \cap \Fi$が
$\Ci(\dS,\, \va)$の真部分集合であることに対応する。

意図の二層化により、「できない」にはもう一種類あることも書ける。
$\Di(\dV, \vs)$が貧しい——やりたい感じ（$\dV$）はあるが、
それを自分の語彙で言い当てる$\dS$が見つからない——のは、
$\Si$の解像度の問題すなわち$\Phii$の問題である。
「頭の中では鳴っているが指が追いつかない」は$\Fi$の問題、
「何かが違うのだが何が違うのか言えない」は$\Phii$の問題であり、
書籍版はこの2つを区別できなかった。

% ------------------------------------------------------------
\subsection{意図の非定常性}\label{subsec:intent-nonstationary}
% ------------------------------------------------------------

$\dV{}^{(t)}$は定数ではない。

ステップ5で$\va^{(t+1)}$を得た後、
再びステップ1に戻って自作を鑑賞すると、
$\vx_i^{(t+1)} = \Pi_i(\va^{(t+1)})$が新たに得られる。
このとき、当初の意図$\dV{}^{(t)}$とは異なる
$\dV{}^{(t+1)}$が生じることがある。
「やってみたら思ったのと違った」
「やってみたらこっちの方が良かった」
という経験がこれに対応する。

さらに、第\ref{sec:dynamics}節で議論する
$\Phii$と$\Psii$の変容がここに接続する。
創作中に$\Phii$が変容すれば$\Si$の地形が変わり、
$\Psii$が変容すれば$\Vi$の地形が変わる。
$\Vi$の地形が変われば、
同じ$\vx_i^{(t)}$に対する$\dV{}^{(t)}$の意味自体が変わる。
追いかけている標的が、追いかけている最中に動くというのが常である。

目的関数自体が探索中に変形するため、創作プロセスは最適化として定式化できない。
候補集合モデルではこの主張がより明確になる。
仮にステップ3・4で何らかの最適な選択基準があったとしても、
$\dV$自体が時変であるため、
プロセス全体を一つの最適化問題として記述することはできない。
これは障害ではなく、C4Pの立場からすれば歓迎すべき構造である。
作り続けるしかない、という結論がここでも導かれる。

% ============================================================
\section{状態量と経路量}\label{sec:state-path}
% ============================================================

\diffstatus{新設。本稿の中心的な再定義であり、
書籍版で「マキシマイズのゲームか否か」と書かれていた区別を、量の型の区別として厳密化する。}

創作ループの反復は軌跡
\begin{equation}
  \traj \;=\; \bigl(\va^{(t)},\, \dV{}^{(t)},\, \dS{}^{(t)},\, \vu^{(t)}\bigr)_{t=1}^{\tau}
\end{equation}
を生成する（$\tau$は停止時刻、第\ref{sec:completion}節）。
このとき：

\begin{itemize}[leftmargin=2em]
  \item \textbf{状態量}：ある時点の状態$\va$（または1ステップ）だけで値が決まる量。
    \[
      Y(\va), \qquad \norm{\bm{\epsilon}_i^{(t)}}, \qquad \Fi \text{の広さ}
    \]
    現在値が存在し、時点ごとに評価でき、大小比較と最適化になじむ。
  \item \textbf{経路量}：軌跡全体$\traj$を引数にとる汎関数$P[\traj]$としてしか定義できない量。
    \[
      \dV \text{の方向の履歴}, \qquad \sigma_i \text{（軌跡からの推定対象として）}, \qquad
      \Phii, \Psii \text{の変容の軌跡}
    \]
    現在値が存在しない。蓄積とともに定まっていき、
    軌跡が短いうちは不確定である。
\end{itemize}

\textbf{C4Qが追う量はすべて状態量であり、C4Pが追う量はすべて経路量である。}
これが本稿におけるC4Q/C4Pの定義である。

この型の違いから、書籍版で個別に論じられていた性質が系として出てくる。

\begin{enumerate}[leftmargin=2em]
  \item \textbf{最適化可能性の非対称}：
    状態量はスカラーの現在値を持つため最大化・最小化の対象になる。
    経路量には最大化すべき現在値がそもそも無い。
    「$\dV$の方向を最大化するとは言えない」（書籍版）は、
    型エラーだから言えないのである。
  \item \textbf{「作り続けるしかない」の定式化}：
    経路量は軌跡が伸びるほど定まる。
    1作品では$\dV$の方向の傾向は不確定だが、蓄積すれば推定が安定する
    （第\ref{sec:trajectory}節）。
  \item \textbf{異時点比較問題の解消}：
    仮定A2により異時点の$\Vi$は同一視できないが、
    経路量は履歴の\textbf{保持}であって異時点値の\textbf{比較}を要求しない。
    「パーソナリティは固定された方向ではなく変容の軌跡そのものである」
    という書籍版の逃げ道は、本稿では定義そのものになる。
  \item \textbf{事後性}：
    経路量は事後的にしか観測されない。
    「作ってみて初めてわかる」は経路量の定義に含まれている。
\end{enumerate}

なお、書籍版の「行為は一つだが変化は2つの量に分解できる」
（第\ref{sec:dynamics}節に収録）は、本稿の言葉では
「同じ軌跡$\traj$から状態量も経路量も読み出せる」と言い直される。
C4Q/C4Pは行為の分類ではなく、\textbf{軌跡から何を読み出すかの分類}である。

% ============================================================
\section{C4QとC4Pの再定義}\label{sec:redefine}
% ============================================================

\diffstatus{改訂。社会的クオリティの定義は変更なし。
技術（実現精度）の節は、ノルム仮定と$\Pi_i$不変性の注記を仮定A2・A3への参照に置き換え、
実現誤差の二分解（記述の失敗／実行の失敗）を追記した。
パーソナリティの構成要素は3つから4つに増えた（停止パターンの追加）。}

候補集合モデルの枠組みを用いて、
クオリティとパーソナリティをそれぞれ構造的に独立な量として定義する。

% ------------------------------------------------------------
\subsection{社会的クオリティ}\label{subsec:social-quality}
% ------------------------------------------------------------

C4Qが言う「客観的な凄さ」は、個人の中で完結する量ではなく、
他者の集合にわたる反応の集計として定義される。
\textbf{総売れ度}
\begin{equation}\label{eq:social-quality}
  Y(\va) = \int_{\I} f\!\bigl((\Psii \circ \Phii)(\va)\bigr)\, d\mu(i)
\end{equation}
がこれに対応する。$\I$は個人の集合、$\mu$は人口分布、
$f : \Vi \to \R$は個人の反応$\vx_i \in \Vi$をスカラーに落とす関数である
（嫌いなものは買わないので、
負の反応をゼロに切る役割を持つ）。

\textbf{C4Qが追求する「クオリティ」とは$Y(\va)$のことである。}
スカラー量であり、マキシマイズのゲームの対象である。
高いに越したことはない。堂々と最大化すればよい。

ただし$Y(\va)$は全個人の$\Pi_i$と人口分布$\mu$を知らなければ計算できない。
個人の立場からは$Y(\va)$は直接観測不可能な量であり、
実際のC4Q的行動は$Y(\va)$の代理指標
（資産額、再生数、売上、フォロワー数など）の最大化として行われる。
本稿のモデルは$Y(\va)$という量が定義できるということだけを主張しており、
その推定方法については関知しない。

% ------------------------------------------------------------
\subsection{技術：実現精度}\label{subsec:technique}
% ------------------------------------------------------------

候補集合モデルにおいて、
創作の1ステップでは選好的意図$\dV{}^{(t)}$に対して
$\dS{}^{(t)}$と$\vu^{(t)}$を選び、$\va^{(t+1)} = T(\va^{(t)},\, \vu^{(t)})$と更新する。
このとき、意図と結果のズレ
\begin{equation}\label{eq:realization}
  \bm{\epsilon}_i^{(t)}
  \;=\;
  \bigl(\Pi_i(\va^{(t+1)}) - \Pi_i(\va^{(t)})\bigr) - \dV{}^{(t)}
\end{equation}
が\textbf{実現誤差}である。
（式\eqref{eq:realization}のノルムと$\Pi_i$の1ステップ不変性については仮定A2・A3を参照。）

$\lVert\epsilon\rVert$が小さいほど、
意図した方向に意図した分だけ$\va$を動かせている。
これが「上手い」ということの定式化であり、
$\Fi$の広さと経験に依存する。

技術の向上とは$\Fi$を拡げ、
$\lVert\epsilon\rVert$を小さくすることであり、
これはマキシマイズ（正確にはミニマイズ）のゲームである。

技術：実現精度と書いたので技巧の話に聞こえやすいが、対象はあらゆる意図である。
ベンチプレス100\,kgを挙げたいとして、40\,kgしか挙げられない状態は
$\lVert\epsilon\rVert$が大きい状態であり、
100\,kg挙げられるようになったら小さくなる。

意図の二層化により、$\bm{\epsilon}$も2つに分解できる。
選んだ$\dS$が実は$\dV$を実現していなかった\textbf{記述の失敗}（耳・目の問題）と、
選んだ$\vu$が$\dS$を実現していなかった\textbf{実行の失敗}（腕の問題）である。
「ミックスで散々いじった結果、いじる前が一番よかった」は前者、
「狙った音程に当たらない」は後者である。

\medskip

技術（実現精度）と社会的クオリティ$Y(\va)$は相関するが別の量である。
技術が高い（$\lVert\epsilon\rVert$が小さい）ことは
$Y(\va)$が高いことの必要条件ではない。
偶然の産物が高い$Y$を持つこともあるし、
技術的に完璧な作品が誰にも響かないこともある。
ただし一般的な傾向として、
意図通りに作れる人間のほうが$Y$を高めやすい。

% ------------------------------------------------------------
\subsection{パーソナリティの再定義}\label{subsec:personality-redef}
% ------------------------------------------------------------

パーソナリティを以下の\textbf{4つの経路量}に分解する
（書籍版の3構成要素に、停止パターンを加えた）。

\begin{enumerate}[leftmargin=2em]
  \item \textbf{選好的意図の履歴}：
    創作プロセスにわたる$\dV{}^{(t)}$の方向の系列。
    「自分はいつもこっちの方向に行きたがる」という傾向。
  \item \textbf{選択カーネル}$\sigma_i$：
    同じ意図に対して$\Di$・$\Ci$から
    どれを選ぶかの傾向。
    「暗くしたい」ときにコードを変える人と
    リバーブを深くする人は、
    $\dV$の方向は同じでも選択が異なる。
  \item \textbf{$\Phii$と$\Psii$の変容の軌跡}：
    創作を通じて$\Phii$（何を知覚できるか）と
    $\Psii$（何を好むか）自体が変わる。
    この変容の履歴。
  \item \textbf{停止パターン}：どこで「完成」と宣言するかの傾向
    （第\ref{sec:completion}節で導入）。
    「詰め切る人」と「勢いを残して止める人」は、
    $\dV$も$\sigma_i$も同じでも異なるパーソナリティである。
\end{enumerate}

パーソナリティは\textbf{プロセスにわたる蓄積}として定義される。
いずれも経路量であり、現在値を持たず、蓄積によってのみ鮮明になる。
$Y(\va)$とは独立である。

% ------------------------------------------------------------
\subsection{C4QとC4Pの分離}\label{subsec:separation}
% ------------------------------------------------------------

以上を踏まえて、C4QとC4Pを定義する。

\paragraph{C4Q（クオリティのための創作）}
\begin{itemize}[leftmargin=2em]
  \item 技術の向上：$\Fi$を拡げ、$\lVert\epsilon\rVert$を小さくする
  \item 社会的クオリティの追求：$Y(\va)$を高める
\end{itemize}
どちらも状態量の改善であり、マキシマイズ（ないしミニマイズ）のゲームである。

\paragraph{C4P（パーソナリティのための創作）}
\begin{itemize}[leftmargin=2em]
  \item $\dV$の方向の蓄積を通じて、自身の選好構造を知る
  \item 選択カーネル$\sigma_i$の蓄積を通じて、自身の行動傾向を知る
  \item $\Phii$と$\Psii$の変容を通じて、新たな自己を発見する
\end{itemize}
マキシマイズではない。
$\dV$の方向を「最大化」するとは言えない（型エラーである。第\ref{sec:state-path}節）。
蓄積し、知り、変容するプロセスである。

\medskip

クオリティは個人レベルのモデルの外に配置される：
社会的クオリティ$Y(\va)$は他者の集合にわたる集計量であり、
個人の内部には存在しない。
個人の内部に残るのは技術（$\Fi$の広さと実現精度$\lVert\epsilon\rVert$）
のみであり、これは「クオリティ」ではなく「技術」と呼ぶ。

ただし、$\va$という共通の変数を通じて実質的に結合している。
$Y(\va)$を追って$\va$を動かすと、
自分の$\dV$との関係が変わる。
市場に合わせた結果「これ、おれが作りたかったものと違うな」
となる現象が、まさにこの結合である。

\textbf{C4QとC4Pは定義上は独立だが、
$\va$を介して緊張関係にありうる。}
この緊張関係の存在こそが、
C4Q/C4Pの区別に実質的な意味を与えている。

% ============================================================
\section{鑑賞と創作の関係}\label{sec:equivalence}
% ============================================================

\diffstatus{記号の更新のみ（$\bm{\delta}_i \to \dV$、選択の二段階化に伴う字句修正）。議論は書籍版のままである。}

創作と鑑賞の関係を、
候補集合モデルの言葉で記述する。

\begin{itemize}[leftmargin=2em]
  \item \textbf{鑑賞}：
    他者が生成した$\va \in \A$に対して$\Pi_i(\va)$を構成する行為。
    パーソナリティは$\Phii$と$\Psii$に宿る。
    $\Pi_i$の順方向のみが用いられる。

  \item \textbf{創作}：
    $\Vi$上の意図$\dV$に基づき、
    $\Di$から$\dS$を、$\Ci \cap \Fi$から$\vu$を選択して
    $\A$上の状態を更新する行為。
    パーソナリティは$\Phii$と$\Psii$に加えて、
    \textbf{選択カーネル$\sigma_i$}にも宿る。
    順方向と逆方向の両方が交互に用いられる。
\end{itemize}

創作は鑑賞を内包する（ステップ1で自作を鑑賞する）。
鑑賞はそれ単独で完結したパーソナリティ発現であり、
創作は鑑賞を部品として含むより大きなプロセスである。

「鑑賞は創作の部分操作」という言い方は、
鑑賞を創作より下に置くように聞こえるかもしれないが、そういう意図はない。
鑑賞をしない創作というものは（このモデルでは）存在しない。
自作を一切鑑賞せずに創作を進めることは、
ステップ1を欠いた候補集合モデルであり、
$\dV$が生成されないため創作プロセスが定義されない。

% ============================================================
\section{軌跡としてのパーソナリティ}\label{sec:trajectory}
% ============================================================

\diffstatus{記号の更新のみ（$\bm{\delta}_i \to \dV$）。
本節の内容は第\ref{sec:state-path}節の「経路量」の具体例として読める。議論は書籍版のままである。}

軌跡モデルは、
候補集合モデルの帰結として自然に再記述される。

候補集合モデルのステップ1--5の反復により、
$\A$上の状態の系列
\begin{equation}
  \bigl\{\va^{(t)} \in \A \bigr\}_{t=1}^{T}
\end{equation}
が生成される。各$\vu^{(t)}$は
$\va^{(t)}$を$\va^{(t+1)} = T(\va^{(t)},\, \vu^{(t)})$へ遷移させる操作であり、
$\Ci(\dS{}^{(t)},\, \va^{(t)}) \cap \Fi$からの選択の結果である。

軌跡を$\A$上のみで記述する代わりに、
三層モデルの構造を踏まえれば、
同時に$\Vi$上の系列
\begin{equation}
  \bigl\{\vx_i^{(t)} = \Pi_i(\va^{(t)}) \in \Vi \bigr\}_{t=1}^{T}
\end{equation}
と、選好的意図の系列
\begin{equation}
  \bigl\{\dV{}^{(t)} \in \Vi \bigr\}_{t=1}^{T}
\end{equation}
も追跡できる。
パーソナリティは$\A$上の軌跡だけでなく、
$\dV$の方向の推移と
候補集合からの選択パターンの推移にも宿る。

$\dV$の方向の履歴によってパーソナリティを推定する。
1作品（短い$\dV$の系列）では方向の傾向は不確定だが、
多数の作品にわたって$\dV$を蓄積すれば、
選好構造の推定は安定する。
パーソナリティの鮮明さは蓄積に依存し、
社会的クオリティ$Y(\va)$とは独立である。

% ============================================================
\section{完成の定式化：停止という選択}\label{sec:completion}
% ============================================================

\diffstatus{新設。本文「完成宣言」の定式化。}

創作ループには自然な停止条件が無い。
すなわち、一般に
\begin{equation}
  \inf_{\vu \in \Fi} \;
  \norm{\,\Pi_i\bigl(T(\va, \vu)\bigr) - \Pi_i(\va) - \dV\,}
  \;>\; 0
  \quad\text{かつ}\quad
  \dV \neq \bm{0}
\end{equation}
が持続する。改善の余地は消えないし、意図も尽きない
（$\dV$は非定常であり、1つ潰せば次が生える）。
したがって\textbf{完成は収束ではない}。
$\norm{\dV} \to 0$を待っていたら作品は永遠に手元に残る。

完成を、創作者が選ぶ\textbf{停止時刻}$\tau$として定義する。
$\tau$はモデルの内部からは導出されない。
「これ以上やらない」という宣言であり、
候補集合からの選択と同格の、もう一つの選択である。

\paragraph{手離しの不可逆化}
本文の「手離しシステム」（物理メディアに残す、人に渡す、締切を切る）は、
停止の\textbf{コミットメント装置}として書ける。
時刻$\tau$以降、その作品に対する実行可能領域を
\begin{equation}
  \Fi^{(t)}(\va) = \emptyset \qquad (t > \tau)
\end{equation}
へ外部の仕組みで固定する操作である。
意志で止まれないなら、$\Fi$の側を外から閉じればよい。
作品のリリースとは、$T$の定義域から自作を除外する手続きである。

\paragraph{内生的停止と外生的停止}
締切による打ち切り（Q\&AのQ11）は外生的停止であり、
完成宣言は内生的停止である。
モデル上の区別は「$\tau$を選んだのが本人の選択か、外部制約か」にある。
そして内生的停止の置き方——
$\dV$との残距離がどれくらいなら手を離せるか——は個人に固有の傾向であり、
第\ref{subsec:personality-redef}節の第4の構成要素（停止パターン）としてパーソナリティに宿る。

% ============================================================
\section{凡庸さと到達可能集合}\label{sec:mediocrity}
% ============================================================

\diffstatus{新設。本文「脱凡庸」の定式化。}

状態$\va$から1ステップで到達できる状態の集合を
\begin{equation}
  \Ri(\va)
  \;=\;
  \bigl\{\, T(\va, \vu) \;\bigm|\; \vu \in \Fi \,\bigr\}
  \;\subset\; \A
\end{equation}
と定義し、\textbf{到達可能集合}と呼ぶ。
$\Fi$が広いほど$\Ri(\va)$は大きい。

\paragraph{同質化のメカニズム}
技術が乏しい段階では、$\Fi$は道具が用意した定型操作
（プリセット、テンプレート、コード進行パターン）の周辺に集中する。
定型は万人に共有されているから、
初学者$i, j$の間で$\Fi \approx \mathcal{F}_j$、
したがって$\Ri(\va) \approx \mathcal{R}_j(\va)$となる。
選べる行き先がほぼ同じなら、
$\sigma_i \neq \sigma_j$（選択の傾向は違う）であっても、
生成される軌跡は$\A$上で互いに似る。
これが\textbf{凡庸さ}——へぼさが類型化したダメさに収束する現象——の構造である。

\paragraph{「技術はパーソナリティの出力経路を増やす」}
$\Fi$の拡大は$\Ri(\va)$を拡げ、
$\sigma_i$の個人差が$\A$上の差として発現する余地を作る。
選択肢が3つしかなければどんな個性も3通りにしか出ないが、
選択肢が1万あれば選択の癖はそのまま軌跡の差になる。
技術向上（C4Qレイヤー1）がC4Pに資する、という本文の主張の構造がこれである。

\paragraph{比較不可能性との整合}
凡庸さの判定は$\A$上（物理状態の近さ）で行われることに注意する。
$\Vi$上の優劣は依然として定義されない。
「作品が似ている」ことと「パーソナリティに優劣がある」ことは別であり、
凡庸さの概念は第\ref{sec:incomparability}節の比較不可能性と矛盾しない。
凡庸とは、パーソナリティが劣っている状態ではなく、
パーソナリティが\textbf{出力に現れない}状態である。

% ============================================================
\section{他者との相互作用：観測の交換}\label{sec:exchange}
% ============================================================

\diffstatus{新設。本文「馴れ合いしようや」の定式化。
書籍版のモデルは個人$i$で閉じており、他者は$Y(\va)$の中の積分変数としてしか現れなかった。}

個人$j$が$i$の作品$\va_i$を鑑賞し、その反応を$i$に報告することを
\begin{equation}
  r_{j \to i} \;=\; g_j\bigl(\Pi_j(\va_i)\bigr)
\end{equation}
と書き、\textbf{観測の交換}と呼ぶ。
$g_j$は$j$の内的反応$\vx_j \in \Vj$を伝達可能な形式（感想・リアクション）に落とす関数である。

$r_{j \to i}$には、受け手$i$にとって2つの異なる使い途がある。

\begin{enumerate}[leftmargin=2em]
  \item \textbf{$Y$の代理指標として}（C4Q的使用）：
    $r_{j \to i}$を$Y(\va_i)$のサンプルとみなし、集計・最大化の手がかりにする。
  \item \textbf{$\Phii, \Psii$の更新入力として}（C4P的使用）：
    「そんな聴き方があるのか」——
    自分の$\Si$に無かった次元の存在を$r_{j \to i}$から知り、
    $\Phii$の変容（第\ref{sec:dynamics}節）の入力にする。
    他者の$\Phij$を借りて自分の知覚の語彙を増やす操作である。
\end{enumerate}

この分離により、本文の馴れ合い3分類が判別可能になる。

\medskip

{\small
\begin{tabular}{@{}lp{52mm}l@{}}
\toprule
分類 & モデル上の記述 & 判定 \\
\midrule
相互ブースト & $r$を協定で歪め、$Y$の代理指標を偽装する & C4Qの不正 \\
社交 & $r$の交換が創作と独立に営まれる & モデル外 \\
C4Pの交換 & $r$を使途2（$\Phii, \Psii$の更新）にのみ用いる & C4P \\
\bottomrule
\end{tabular}
}

\medskip

相互ブーストが不正なのは、代理指標と$Y$の対応を壊すからである
（指標は上がるが$Y$は上がっていない）。
C4Pの交換はそもそも$Y$を参照しないため、歪めるべき指標が無い。
「馴れ合い」への批判は使途1の文法でなされ、
本文が擁護するのは使途2である。
両者は$r_{j \to i}$という同じ信号の異なる使い方であり、
外形からは区別しにくいが、モデル上は明確に別の操作である。

なお、複数人の$\dV$を一つに統合する操作は引き続き定義しない
（Q\&AのQ24と同じ立場）。
交換できるのは観測であって、意図ではない。

% ============================================================
\section{運の所在}\label{sec:luck}
% ============================================================

\diffstatus{新設。本文「不運な人」「アウトオブスコープ」の定式化。}

本モデルの外生変数——本人の選択が及ばず、初期条件として与えられるもの——を列挙する。

\begin{itemize}[leftmargin=2em]
  \item $\Fi^{(0)}$：初期実行可能領域。身体・機材・可処分時間・経済基盤。
  \item 刺激の到達分布：どんな作品・概念に出会える環境か。
    $\Phii, \Psii$の変容（第\ref{sec:dynamics}節）の入力は環境から来る。
  \item $\mu$：人口分布。$Y(\va)$の定義そのものに入っている。
\end{itemize}

状態量はこれらに直接感染する。
$Y(\va)$は$\mu$の関数だから、同じ$\va$でも時代と場所で$Y$は変わる。
$\Fi$の拡大速度は$\Fi^{(0)}$と可処分時間に依存する。
\textbf{C4Qのゲームは初期条件の抽選結果を織り込んだゲームである。}
尺度を導入すれば、運良く「優」側に入れた者が利する——本文の主張のとおりである。

一方、経路量の\textbf{定義}は外生変数を含まない。
$\dV$の履歴、$\sigma_i$、変容の軌跡、停止パターンのどれも、
$\mu$にも$\Fi^{(0)}$にも言及せずに定義されている。
「パーソナリティだけが運に作用されず全員に平等」の正確な意味はこれである：
\textbf{定義可能性が環境に依存しない。}
軌跡が1歩でも伸びれば経路量は蓄積し始め、
その蓄積は誰の環境とも比較されない。

ただし正直に書くべき留保が2つある。
蓄積の\textbf{速度}は環境に依存する（ループが回らなければ蓄積は止まる）。
そして、ループを回すための前提条件（機材・身体・時間）を欠く人に対して
本モデルは何も提供しない。
C4Pは態度の理論であって社会福祉の理論ではない——
本文「アウトオブスコープ」の宣言は、モデル上は
「外生変数はモデルの外にある」という一行に対応する。
平等なのは定義であって、境遇ではない。

% ============================================================
\section{写像の変容と候補集合モデルの接続}\label{sec:dynamics}
% ============================================================

\diffstatus{改訂。書籍版の内容（$\Fi$の拡大が$\Phii$の変容を引き起こす／行為の不可分性）はそのまま維持し、
冒頭に更新規則$U_\Phi, U_\Psi$による抽象的な記述を追加した。
変容の入力として観測の交換$r_{j \to i}$（第\ref{sec:exchange}節）が加わった。}

$\Phii$と$\Psii$の時間発展を、抽象的な更新規則
\begin{equation}
  \Phii^{(t+1)} = U_\Phi\bigl(\Phii^{(t)},\, e^{(t)}\bigr),
  \qquad
  \Psii^{(t+1)} = U_\Psi\bigl(\Psii^{(t)},\, e^{(t)}\bigr)
\end{equation}
と書く。$e^{(t)}$は時刻$t$の\textbf{経験}であり、
自作の制作、他者作品の鑑賞、観測の交換$r_{j \to i}$のいずれでもよい。
$U_\Phi, U_\Psi$の関数形は与えない（仮定A5）。
形式を与えたのは、変容の入力経路を明示するためだけである。
鑑賞専門の人にもパーソナリティの変容が起きること（Q\&AのQ13）、
スランプや制作不能期にも地形が動き続けること（Q25・Q28）は、
$e$が他者作品のみの場合として統一的に記述される。

\subsection{$\Fi$の拡大が$\Phii$の変容を引き起こす}

$\Fi$が拡がる（技術が向上する）と、
$\A$上の新たな領域を探索可能になる。
その結果、$\Phii$が接触する刺激の範囲が広がり、
$\Si$の次元追加が生じうる。

\textbf{例}：ピアニストがショパンのエチュードを弾けるようになった結果、
以前は弾けなかったフレーズを弾くことで
「この和声の移り変わりにこんな感覚があったのか」と発見する。
$\Fi$の拡大が$\Phii$の変容を誘発した例である。

これはC4Q的行為（$\Fi$の拡大）が
C4P的帰結（$\Phii$の変容による自己発見）を
副産物として生むメカニズムの記述である。

\subsection{行為の不可分性と量の分離可能性}

技術を極める過程で自己発見が起きることは珍しくなく、
むしろ常態である。
書道家が何千回も同じ字を書く中で、
技術が上がると同時に「自分の線」が見えてくる。

このとき、\textbf{行為は一つ}（同じ字を書く反復）だが、
\textbf{行為が生む変化は2つの量に分解できる}：
\begin{itemize}[leftmargin=2em]
  \item $\Fi$が拡がっている（筆の制御精度が上がる）
  \item $\dV$の履歴が蓄積している
    （自分がどの線を選ぶ傾向にあるかが見えてくる）
\end{itemize}

行為が不可分であっても、行為が生む変化は分離可能である。
C4Q/C4Pは行為の分類ではなく、\textbf{何を追っているかの分類}である
（第\ref{sec:state-path}節の言葉では、状態量を読むか経路量を読むかの分類である）。

C4Q的行為（$\Fi$の拡大）がC4P的帰結（$\Phii$の変容）を
副産物として生むことは珍しくなく、むしろ常態である。
「過去の自分を超える」という体験がC4QとC4Pのどちらに該当するかは
Q\&AのQ5を参照されたい。

% ============================================================
\section{比較不可能性について}\label{sec:incomparability}
% ============================================================

\diffstatus{ほぼ変更なし。末尾に、同一個人の異時点間の問題が経路量の定義で解消されたことと、
凡庸さとの整合についての2段落を追加した。}

過去の検討で
「個人間の$\Si$, $\Vi$の比較が非自明である」
「共通の潜在空間への埋め込みが必要になるが、
これは元の問題（外生的な共通空間）を再導入するリスクがある」
と述べた。

これを\textbf{課題ではなくモデルの帰結}として受け入れる。

異なる個人$i, j$の$\Vi$と$\Vj$は一般に異なる空間であり、
$\dV_i$と$\dV_j$を直接比較することはできない。
パーソナリティの比較が非自明であることは、
「パーソナリティに優劣はない」というC4Pの主張の
数理的な帰結である。
比較するための共通空間を導入した時点で、
よさの定義権が個人から外部に移るため、
三層モデルの趣旨に反する。
比較不可能性は欠陥ではなく、本モデルでは設計通りであるとする。

本稿での追加は2点である。
第一に、仮定A2により同一個人の異時点間にも同じ構造があることが明示され、
経路量による定義（第\ref{sec:state-path}節）がその両方への統一的な回答になった。
第二に、凡庸さ（第\ref{sec:mediocrity}節）は$\A$上の近さの話であり、
$\Vi$上の比較を要求しないため、本節と矛盾しない。

% ============================================================
\section{本モデルの射程}\label{sec:scope}
% ============================================================

\diffstatus{改訂。書籍版で保留されていた論点のうち、
意図の所在（→第\ref{subsec:intent}節）、選択分布（→第\ref{subsec:kernel}節）、
異時点比較（→第\ref{sec:state-path}節）は本稿で回収した。
残る保留を整理し、楽しみ力・記憶についての項目を追加した。}

構造によって態度を示すことが目的なので、
$\Phii$と$\Psii$の具体的な関数形は与えていない。
それによって、このモデルを用いて何かを具体的に数値化するなどの
応用先は現状存在しない。
本モデルは態度を表明するための構造であり、
汎用的な解析ツールとして機能させることは意図していない。

残る保留は以下のとおりである。

\begin{itemize}[leftmargin=2em]
  \item $\Phii, \Psii, U_\Phi, U_\Psi, \sigma_i$の関数形は与えない（仮定A5）。
  \item 「楽しさ」は定式化しない。
    C4Pの継続は、経路量そのものに価値を感じられるかに懸かっており、
    本文「楽しみ力」はそれが養成可能な能力だと主張するが、
    「経路に対する選好」を$\Psii$と別立てで導入すると
    モデルが急激に重くなるため、本稿では踏み込まない。
    ただし位置づけだけ書いておく：
    C4Qは状態量に価値を置く態度、
    C4Pは経路量に価値を置く態度であり、
    楽しみ力とは後者の価値を実際に感じ取る能力のことである。
  \item 記憶・忘却・一回性の感動（本文「あれなんだっけ？」「スペシャル（for me）」）は扱わない。
    軌跡$\traj$は完全な記録として書かれているが、
    実際の創作者は自分の軌跡を忘れ、歪めて想起する。
    忘却込みの軌跡モデルは将来の課題として面白いが、本稿の射程外である。
  \item 倫理と経済（本文「人間と作品」の棚上げ部分）は引き続き扱わない。
\end{itemize}

% ============================================================
\section{まとめ}
% ============================================================

\diffstatus{改訂。書籍版のまとめを、本稿の新設概念を含む形に更新した。}

本稿で構築したモデルの骨格を整理する。

\begin{itemize}[leftmargin=2em]
  \item 創作物の物理的状態が張る属性空間$\A$、個人$i$の知覚・認知が畳み込む意味空間$\Si$、好悪の座標系である選好潜在空間$\Vi$の三層を射影$\Pi_i = \Psii \circ \Phii$で接続した。
  \item 意図を選好的意図$\dV \in \Vi$と記述的意図$\dS \in \Si$に二層化し、創作を二段階の逆問題として定式化した。選択点は2箇所あり、どちらの選択にもパーソナリティが発現する。技術の不足も「言えない（$\Phii$の問題）」と「できない（$\Fi$の問題）」に分離される。
  \item C4QとC4Pの区別を「状態量／経路量」という量の型の区別として定義した。C4Qは状態の関数（$Y(\va)$、$\lVert\epsilon\rVert$、$\Fi$の広さ）を改善するゲームであり、C4Pは軌跡の汎関数（$\dV$の履歴、選択カーネル$\sigma_i$、$\Phii$・$\Psii$の変容の軌跡、停止パターン）の蓄積である。
  \item 経路量に現在値は無く、最適化は型エラーであり、事後的にしか観測されない。「作ってみて初めてわかる」はこの型の性質である。
  \item C4Qが追求する「クオリティ」は個人の内部に存在しない。他者の集合にわたる反応の集計量$Y(\va)$として社会の側に配置される。個人の内部に残るのは技術であり、クオリティとパーソナリティは構造的に独立な量である。
  \item 鑑賞は創作プロセスの部分操作として位置づけられる。創作は鑑賞を内包し、鑑賞をしない創作はこのモデルでは定義されない。
  \item 完成は収束ではなく、停止時刻$\tau$の選択である。停止パターンはパーソナリティの第4の構成要素である。
  \item 凡庸さは到達可能集合$\Ri(\va)$の個人間の重なりとして記述される。技術向上はパーソナリティの出力経路を増やす。凡庸とはパーソナリティの劣後ではなく、出力に現れない状態である。
  \item 観測の交換$r_{j \to i}$の2つの使途（$Y$の代理指標／$\Phii, \Psii$の更新入力）により、馴れ合いの3分類が判別可能になる。
  \item 運はモデルの外生変数（$\Fi^{(0)}$・刺激の到達分布・$\mu$）に局在する。状態量は運に感染するが、経路量の定義は環境に依存しない。平等なのは定義であって境遇ではない。
\end{itemize}

$\dV$の方向は本人にしかわからないし、事前に知ることもできない。
作ってみて初めてわかる。それがC4Pである——
この結論は書籍版と同じである。変わったのは、
この文が比喩ではなく経路量の定義から出てくるようになったことである。

\clearpage

% ============================================================
\section*{Q\&A}\label{sec:qa}
% ============================================================

\diffstatus{Q1--Q30は書籍版のまま（記号$\bm{\delta}_i \to \dV$と、ステップ番号・参照先のみ本稿に合わせて更新）。
Q31--Q36は本稿で追加した。}

本モデルに対してよくされる質問と、モデル上の回答をまとめる。

\qa{何でこんなことを始めたの？}
考え方を示すのが目的であって、モデルを運用して何かをしようとはしていない。
C4P作文をあてもなく書きまくった結果、
自分が何を主張しているのかを自分で整理する必要が出てきたからである。
「クオリティとパーソナリティは独立」
「客観的クオリティは幻想」
「創作を通じて自分を知る」
といった主張を繰り返しているうちに、
それぞれの主張が互いに整合しているのか、
矛盾していないかが気になりはじめた。
日本語は便利だが曖昧なので、
構造を一回ちゃんと書いて確認したかった。
定式化の過程で矛盾を発見し解消するなど、この整理は本文の議論にもよい影響を与えた。最初は単なる数理ポエムであったが、思った以上に実務的なよい効果があった。

\qa{数式にする意味はあるのか。日本語で十分説明できるのでは}

本モデルの価値は、日本語の議論だけでは曖昧になる区別に
判定基準を与えることにある。加えて、おんなじ質問ばっかされるのに日本語だけでは毎回説明が面倒なので、回答をするためには構造を正確に記述する必要があった。
Q5がその典型例で、「過去の自分を超える」が
$\Fi$の拡大なのか$\dV$の方向変化なのか、
日本語では一括りにされる体験をモデル上で分離できる。
Q4もそうで、「速弾きが好き」と「速弾きが上手い」を
$\Psii$と$\Fi$として別の量に帰属させることで、
日常語の「パーソナリティ」が指しうる複数の意味を整理できる。
この分離が使えるかどうかは読者次第であるが、
少なくとも筆者は使えている。

\qa{自分に都合のいい理屈を作っているだけでは？}

ほぼYES。自分の文章に都合よくなるようにモデルを作っている。
ただし、都合のいい理屈であっても構造が内部的に整合していれば、
少なくとも自分の態度の表明としては機能する。

加えて、構造を書くことで都合の良さに起因する矛盾が検出可能になった。事後的であるが、都合のいい理屈を精密に書くと、都合の悪い部分が浮き上がるという数理モデルの副次的な効能に気がついた。

\qa{「技術だけを高め続けたい」はC4Q的態度だが、それもパーソナリティと呼べるのでは？}

呼べる。ただし、本モデルにおけるパーソナリティの定義とは位相が異なる。

例えば、ギタリストが速弾きの精度を極めることだけに関心がある場合、
この人は$\Fi$の拡大を追求している。
モデル上これはC4Qのレイヤー1（技術）であり、
$\dV$の方向やその蓄積とは独立なパラメータである。

しかし、「速弾きを極めたい」という志向そのものは、
この人の選好構造$\Psii$に由来している。
遅いアルペジオではなく高速フルピッキングに魅力を感じる、
という$\Vi$上の方向は確かにこの人に固有のものであり、
日常語では「個性」「パーソナリティ」と呼ばれうる。

本モデルの用語法では、$\Psii$の構造（何に魅力を感じるか）は
パーソナリティの構成要素であるが、
$\Fi$の拡大行為そのものはパーソナリティの発現ではない。
「速弾きが好きだ」はパーソナリティ、
「速弾きが上手い」は技術。
好きであることと上手いことは別の量であり、
本モデルではこの区別を維持する。

さらに、速弾きを極める過程で
「自分は実は速さそのものではなく、
速いパッセージの中の微妙なダイナミクスの変化に惹かれていたんだな」
と気づくことがある。
これは$\Psii$の変容であり、C4P的帰結である。
$\Fi$の拡大（C4Q的行為）が$\Psii$の変容（C4P的帰結）を
副産物として引き起こした例であり、
行為は一つだが変化は2つの量に分解できる。

\qa{「過去の自分を超えていく＝質を高める」はC4Q？C4P？}

どちらでもありうる。何を超えているかによって決まる。

\medskip

{\footnotesize\begin{tabular}{@{}l@{\hspace{3mm}}l@{\hspace{3mm}}l@{}}
\toprule
何が変わったか & モデル上の対応 & 判定 \\
\midrule
以前できなかった技法ができるようになった & $\Fi$の拡大 & C4Q \\
他人に「良くなった」と言われた & $Y(\va)$の増加 & C4Q \\
意図通りの音が出せるようになった & $\lVert\epsilon\rVert$の減少 & C4Q \\
\midrule
前と違うものを作りたくなった & $\dV$の方向変化 & C4P \\
以前は聴こえなかった要素が聴こえるようになった & $\Phii$の変容 & C4P \\
自分がこれを好きだったと気づいた & $\Psii$の変容 & C4P \\
自分の癖や傾向が見えてきた & $\dV$の蓄積 & C4P \\
\bottomrule
\end{tabular}}

\medskip

「過去の自分を超えていく＝質を高める」という等式自体が、
2つの異なるプロセスを束ねている。
それを分解する道具を提供するのが本モデルの仕事である。

\qa{インプットが大事って言うけど、それはC4Q？C4P？}

場合による。何のためのインプットかで決まる。

\medskip

{\footnotesize\begin{tabular}{@{}l@{\hspace{3mm}}l@{\hspace{3mm}}l@{}}
\toprule
何をしているか & モデル上の対応 & 判定 \\
\midrule
特定の技法を学ぶために教材を摂取 & $\Fi$の拡大 & C4Q \\
\midrule
多くの作品を鑑賞し新たな概念を獲得 & $\Phii$の変容（$\Si$の次元追加） & C4P \\
鑑賞を通じて自分の好みに気づく & $\Psii$の変容 & C4P \\
\bottomrule
\end{tabular}}

\medskip

同じインプットがC4QとC4Pの両方に効くことは珍しくない。

\qa{リファレンス曲に寄せるのはC4Q？C4P？}

場合による。何のために寄せているかで決まる。

\medskip

{\footnotesize\begin{tabular}{@{}l@{\hspace{3mm}}l@{\hspace{3mm}}l@{}}
\toprule
何をしているか & モデル上の対応 & 判定 \\
\midrule
売れている曲に寄せて$Y$を上げたい & $Y(\va)$の最大化 & C4Q \\
あの手法を自分でも再現できるようになりたい & $\Fi$の拡大 & C4Q \\
\midrule
あの曲のここが好きで、自分でもやりたい & $\dV$に従った$\vu$の選択 & C4P \\
寄せる過程で新しい構造に気づいた & $\Phii$の変容 & C4P \\
\bottomrule
\end{tabular}}

\medskip

多くの場合、これらは同時に起きている。

\qa{好きなアーティストに影響を受けているのはC4Q？C4P？}

影響の受け方による。

\medskip

{\footnotesize\begin{tabular}{@{}l@{\hspace{3mm}}l@{\hspace{3mm}}l@{}}
\toprule
何が起きているか & モデル上の対応 & 判定 \\
\midrule
あの手法を使えるようになりたい & $\Fi$の拡大 & C4Q \\
あのアーティストみたいに売れたい & $Y(\va)$の追求 & C4Q \\
\midrule
聴こえなかった要素が聴こえるようになった & $\Phii$の変容 & C4P \\
自分もこういう方向に行きたいと感じた & $\dV$の方向形成 & C4P \\
\bottomrule
\end{tabular}}

\medskip

多くの場合、これらは同時に起きている。

\qa{コンペに出すのはC4Q？C4P？}

基本的にC4Qである。

\medskip

{\footnotesize\begin{tabular}{@{}l@{\hspace{3mm}}l@{\hspace{3mm}}l@{}}
\toprule
何が起きているか & モデル上の対応 & 判定 \\
\midrule
入賞を目指している & $Y(\va)$の最大化ゲームへの参加 & C4Q \\
審査基準に合わせて$\va$を調整 & $Y(\va)$の代理指標の最大化 & C4Q \\
\midrule
仕上げる過程で自分の方向が見えた & $\dV$の蓄積 & C4P（副産物） \\
\bottomrule
\end{tabular}}

\medskip

「コンペがなければ作らなかった」作品から自己発見が起きることは珍しくない。

\qa{何も考えずに手を動かしてたら良いのができた。これはC4Q？C4P？}

行為の時点では判定できない。結果で決まる。

\medskip

{\footnotesize\begin{tabular}{@{}l@{\hspace{3mm}}l@{\hspace{3mm}}l@{}}
\toprule
事後に何が起きたか & モデル上の対応 & 判定 \\
\midrule
他人に「いいね」と言われた & $Y(\va)$が高かった & C4Q的成果 \\
\midrule
「自分はこういうのが好きだったのか」と気づいた & $\dV$の事後的な蓄積 & C4P的成果 \\
\bottomrule
\end{tabular}}

\medskip

$\dV$が不明確なまま$\vu$を選んでいる状態であり、
ステップ2（選好的意図の認知）が曖昧なままプロセスが進行している。
行為が先で意図の認知が後、という順序はモデル上も許容される。

\qa{締め切りがあるから妥協するのはC4Q？C4P？}

どちらでもない。モデルの外の制約条件である。

\medskip

{\footnotesize\begin{tabular}{@{}l@{\hspace{3mm}}l@{\hspace{3mm}}l@{}}
\toprule
何が起きているか & モデル上の対応 & 判定 \\
\midrule
$\lVert\epsilon\rVert$を小さくしきれず打ち切り & プロセスの中断 & C4Q/C4Pの外 \\
\midrule
「これでいいか」と思える & $\dV$との距離が許容範囲 & --- \\
「これは自分の作品ではない」と感じる & $\dV$との距離が大きい & C4P的な不満 \\
\bottomrule
\end{tabular}}

\medskip

打ち切りの判断は時間や予算の制約であり、
C4Q/C4Pのどちらを追っているかとは独立である
（本稿の言葉では外生的停止である。第\ref{sec:completion}節参照）。

\qa{作ったけど公開しないのはC4Q？C4P？}

C4Pとしては問題ない。損失はC4Q側にのみ生じる。

\medskip

{\footnotesize\begin{tabular}{@{}l@{\hspace{3mm}}l@{\hspace{3mm}}l@{}}
\toprule
何が起きているか & モデル上の対応 & 判定 \\
\midrule
自作を鑑賞し意図との距離を感じた & $\dV$の蓄積 & C4Pは機能 \\
$Y(\va)$のフィードバックが得られない & 代理指標の学習が停止 & C4Qに損失 \\
\bottomrule
\end{tabular}}

\medskip

$\dV$の蓄積は公開の有無に依存しない。

\qa{鑑賞専門の人にもパーソナリティはあるの？}

ある。
鑑賞は$\Pi_i$の順方向適用であり、
$\Phii$と$\Psii$にパーソナリティが宿る。
同じ楽曲を聴いて異なる感想を持つこと自体が、
$\Phii \neq \Phij$かつ/または$\Psii \neq \Psij$の帰結であり、
パーソナリティの発現である。
創作しなくてもパーソナリティは発現する。

\qa{このモデル、音楽以外にも使えるの？}

使える。小説、絵画、料理、プロダクトデザイン、なんでもよい。
$\A$上の状態の物理的内容が変わるだけで、
三層射影の構造と候補集合モデルはそのまま適用される。
本稿で例が音楽に偏っているのは筆者が音楽家だからであり、
モデルの制約ではない。

\qa{属性空間が無限次元って言われてもピンとこない}

CDに収録された44.1\,kHz、ステレオ、3分間の音源は
約1500万次元のベクトルである（第\ref{sec:three-layers}節参照）。
これだけで十分にバカでかい。
実際の創作物がこの空間のごく薄い部分多様体上に乗っている、
という直感が持てればそれで十分である。

\qa{売れたいと思うのはダメなことなの？}

ダメではない。
$Y(\va)$を追うことは正当なC4Qであり、
堂々と最大化すればよい。
C4Pは$Y(\va)$の追求とは別のゲームだと言っているだけであって、
$Y(\va)$を追うこと自体を否定していない。

\qa{C4QとC4Pを両立することはできるの？}

できる。というか、どちらかだけやるのは難しいと思う。
$\va$を共有するため緊張関係は生じうるが（第\ref{subsec:separation}節参照）、
多くの創作者は実際に両方を同時に追っている。
ただし、$Y(\va)$の追求が$\dV$の方向を歪めていると感じるなら、
それはC4Q側の圧力がC4P側に干渉しているということであり、そういったものを問題視している。
どちらを優先するかは本人の態度の問題である。

\qa{パーソナリティなんて気にしなくても良い作品は作れるのでは？}

作れる。
$Y(\va)$が高い作品を作ることとパーソナリティの発現は独立である。
C4Qだけで十分な人はそれでよい。
C4Pは全員に推奨しているが強制ではない。

\qa{AIが作った曲にパーソナリティはあるの？}

本モデルではパーソナリティは
$\Phii$、$\Psii$、$\dV$の方向、
候補集合からの選択パターンに宿る。
これらはすべて個人$i$に属する量である。
AIが$\va \in \A$を生成しても、
それを鑑賞する人間のパーソナリティは$\Pi_i(\va)$を通じて発現する。
一方、AI自体が$\Phii$や$\Psii$に相当するものを持つかどうかは、
本モデルの射程外である。

\qa{「AIに曲をつくらせる」はこのモデルでどう扱える？}

候補集合モデルの枠組みで扱える。
AIに指示を出す行為は、候補集合からの選択$\vu$に対応する。
プロンプトやパラメータの選択が$\vu$であり、
AIの出力が$\va' = T(\va,\, \vu)$である。実現誤差$\bm{\epsilon}_i$は通常の創作と同様に伴う。

自分の手で楽器を弾く場合と構造は変わらない。
変わるのは$\Fi$（実行可能領域）の形状である。
AIを道具として使うことで、
自力では到達できなかった$\A$上の領域に到達可能になる。
これは$\Fi$の拡大であり、本文の「サイズはどうしますか」章で論じた
工程の粒度の問題と同じ構造である。

C4P的に重要なのは、AIの出力に対して
「これは自分が意図したものか」を判断するプロセスである。
この判断は$\dV$と$\va'$の距離の評価であり、
人間の$\Psii$を通じて行われる。
AIが生成しても、選択と判断が人間に属する限り、
$\dV$の蓄積は生じる。

\qa{パーソナリティが鮮明じゃない人はどうすればいいの？}

$\dV$の蓄積が足りないだけである。
作り続ける、あるいは鑑賞し続ける。
蓄積すれば方向推定が安定する（第\ref{sec:trajectory}節参照）。
量の問題であり、才能の問題ではない。

\qa{好きなものがコロコロ変わるのはパーソナリティがないということ？}

違う。
$\Psii$の変容が活発なだけであり、
パーソナリティがないのではなく、
変容の軌跡そのものがパーソナリティである。
固定されていることがパーソナリティの条件ではない。

\qa{他人の評価が気になってしまうのはC4Qに毒されているということ？}

$Y(\va)$の代理指標（再生数、いいね数など）を気にしているだけであり、
毒されているかどうかは態度の問題である。
気にすること自体がダメなのではなく、
それが$\dV$の方向を歪めているかどうかが論点である。
代理指標を参考にしつつ$\dV$の方向を保つことは可能であり、
その場合C4QとC4Pは両立している。

\qa{複数人でのプロジェクトではC4Pはどう扱われる？}

候補集合モデルは個人$i$を主語に書かれている。
複数人のプロジェクトでは、メンバーごとに$\dV$が異なる。
$Y(\va)$は共有できる（同じ作品の売上は全員に帰属する）が、
$\dV$の方向は個人に属するため共有できない。

バンドやチームで「方向性の違い」が生じるのは、
メンバー間で$\dV$の方向が乖離している状態である。
$Y(\va)$の最大化では合意できても、
候補集合からどの$\vu$を選ぶかで衝突する。
C4Q的には協調可能だがC4P的には本質的に個人の営みであり、
複数人の$\dV$を一つに統合する操作は本モデルでは定義されない。

\qa{創作できない状態（金がない、時間がない、病気など）の人はC4Pではどう扱われる？}

候補集合モデルのステップ1--5のループが回らない状態に対応する。
ループが停止していれば$\dV$の新たな蓄積は生じず、
パーソナリティの発現プロセスは一時停止する。

ただし、鑑賞は可能である場合が多い。
$\Pi_i$の順方向適用（他者の作品を鑑賞する）は
創作のループとは独立に機能し、
$\Phii$や$\Psii$の変容は鑑賞を通じて起きうる。
創作できない期間に鑑賞を通じて$\Phii$が変容し、
復帰後に「以前とは違うものを作りたくなった」と感じることは
このメカニズムで説明される。

C4Pは「今すぐ作れ」とは言っていない。
蓄積のペースは人それぞれであり、
停止期間があってもパーソナリティが消失するわけではない。
金がない・時間がない・病気であるといった制約は、
Q11（締め切り）と同様にモデルの外の制約条件として扱われる。
C4Pはこれらの制約を解消する手段を持たない。
C4Pは創作に対する態度の理論であって、社会福祉の理論ではない。

\qa{才能がないと思うんですが、それでもC4Pは機能しますか？}

機能する。才能という語が指しうるものをモデル上で分解する。

\medskip

{\footnotesize\begin{tabular}{@{}l@{\hspace{3mm}}l@{\hspace{3mm}}l@{}}
\toprule
何を指しているか & モデル上の対応 & C4Pとの関係 \\
\midrule
技法の習得が早い & $\Fi$の拡大速度が速い & 無関係 \\
作ったものが人に評価される & $Y(\va)$が高い & 無関係 \\
\midrule
独自の感性がある & $\Phii$/$\Psii$が他者と大きく異なる & パーソナリティそのもの \\
\bottomrule
\end{tabular}}

\medskip

最初の2つはC4Qに属する量であり、C4Pの機能とは独立である。
3つ目については、$\Phii$と$\Psii$は全員が異なるのだから、
「独自の感性がない」人間はモデル上存在しない。
才能がないと感じている人は、
$\dV$の蓄積が足りないか、
$Y(\va)$の低さを才能のなさと混同しているかのどちらかである。

\qa{自分の作品を真似された／パクられた場合、パーソナリティは奪われますか？}

奪われない。
真似されるのは$\va$（属性空間上の状態）であり、
$\Pi_i$は複製できない。
ある人が自分の曲を完全にコピーしたとしても、
そのコピーに至る$\dV$の履歴と候補集合からの選択パターンは
コピー元の本人にしか属さない。
表面的に同じ$\va$が存在しても、
それを生成した選択の軌跡は異なる。

逆に言えば、真似する側にもその人固有の$\Pi_i$を通じた鑑賞と選択が発生しており、
「何を真似するか」の選択自体にパーソナリティが宿る。

\qa{スランプで作れなくなった状態はどう扱われますか？}

Q25（外的制約で作れない）とは異なり、スランプは内的な状態である。
モデル上はいくつかの記述が可能である。
$\dV$が生成されない（何をしたいかわからない）場合、
ステップ2（選好的意図の認知）が機能していない。
$\Phii$や$\Psii$が変容の途中にあり、
以前の$\dV$の方向が現在の選好と合わなくなっている場合もある。

いずれにせよ、$\Phii$と$\Psii$は鑑賞を通じて変容し続けるため、
作れない期間にも地形は動いている。
スランプの「出口」は、
変容後の選好構造に合致する$\dV$が再び生成される時点である。

\qa{趣味でやってるだけなんですが、C4Pとか大げさでは？}

大げさではない。
候補集合モデルはプロとアマチュアを区別するパラメータを持たない。
週末に絵を描く人も、毎日曲を作る人も、
候補集合からの選択という構造は同じである。
$\dV$の蓄積速度は頻度に依存するが、
蓄積の構造自体は変わらない。
「趣味だからパーソナリティは発現しない」とはならない。

むしろ、外部からの要請（$Y(\va)$の最大化圧力）が弱い分、
$\dV$の方向が歪みにくいという意味で、
趣味の創作はC4P的に純度が高い状態であるとも言える。

\qa{結局、何を作ればいいの？}

知らない。このモデルはその問いに答えない。
$\dV$の方向は本人にしかわからないし、
事前に知ることもできない。
作ってみて初めてわかる。それがC4Pとも言える。

% ------------------------------------------------------------
% 以下、本稿で追加したQ&A（Q31--Q36）
% ------------------------------------------------------------

\qa{「完成した」といつ言えばいいの？}

モデルは答えない。
完成は収束ではなく選択であり、$\tau$を導出する式は存在しない（第\ref{sec:completion}節）。
モデルが言えるのは2つだけである。
改善余地がゼロになるのを待つ戦略は停止しないこと。
そして、どこで止めるかの傾向自体がパーソナリティの構成要素であること。
「詰め切ってから離す人」と「勢いを残して離す人」の違いは、
好みや技術の違いと同格の、あなたの一部である。

\qa{初心者の作品が「どれも似てる」のはパーソナリティがないから？}

違う。$\Fi$が定型操作に集中しているため、
到達可能集合$\Ri(\va)$が他の初心者とほぼ重なっているだけである
（第\ref{sec:mediocrity}節）。
選択の傾向$\sigma_i$は最初から個人ごとに異なるが、
選択肢が3つしかなければ個性は3通りにしか出力されない。
技術が上がると$\Ri$が拡がり、同じ$\sigma_i$が違う軌跡として現れ始める。
パーソナリティが無いのではなく、出力の解像度が足りない。

\qa{仲間内で褒め合ってるだけって言われた}

褒め合いの中身による（第\ref{sec:exchange}節）。
互いの代理指標（いいね数・再生数）を協定で吊り上げているなら相互ブーストであり、
C4Qの文法においても不正である。
互いの反応$r_{j \to i}$を交換して
「そんな聴き方があるのか」を持ち帰っているなら、
それは$\Phii, \Psii$の更新入力の交換であり、C4Pとして正当に機能している。
外形は似るがモデル上は別の操作である。
批判者が前者の文法で後者を撃っているなら、その批判は当たっていない。

\qa{「何が違うのか言えないけど何かが違う」状態はどう扱われる？}

記述候補集合$\Di(\dV, \vs)$が貧しい状態である（第\ref{subsec:intent}節）。
$\dV$（もっとこうしたい）は立っているが、
それを自分の語彙$\Si$で言い当てる$\dS$が見つからない。
これは腕（$\Fi$）ではなく耳・目（$\Phii$の解像度）の問題であり、
処方も異なる：練習ではなく、鑑賞と概念の獲得が効く。
逆に「言えるけどできない」は$\Fi$の問題で、こちらは練習が効く。
書籍版はこの2つを区別できなかったが、意図の二層化で分離された。

\qa{環境に恵まれてる人が有利なのは結局変わらなくない？}

変わらない。モデルはそこを糊塗しない（第\ref{sec:luck}節）。
$Y(\va)$も$\Fi$の拡大速度も初期条件の抽選に感染しており、
C4Qのゲームは運を織り込んだゲームである。
モデルが言えるのは、経路量の\textbf{定義}が環境を参照しないことまでである。
あなたの$\dV$の履歴は誰の環境とも比較されずに蓄積する。
ただし蓄積の速度は環境に依存するし、
ループを回す前提条件を欠く人にこのモデルは何も配れない。
平等なのは定義であって境遇ではない。

\qa{書籍版と結論が変わったところはあるの？}

主張の結論は変わっていない。変わったのは導出の形である。
書籍版で「比喩」「保留」「脚注」だったものが定義と系に昇格した。
たとえば「作り続けるしかない」は書籍版では激励に聞こえたが、
本稿では経路量の定義（現在値が無く、蓄積によってのみ定まる）からの帰結である。
「パーソナリティに優劣はない」に加えて
「凡庸さはパーソナリティの劣後ではない」が言えるようになったのは、
凡庸さを$\A$上の近さとして定義し、$\Vi$上の比較と切り離したからである。
数式が増えて主張が過激になるのではなく、
同じ主張がより少ない言い訳で立つようになるのが改訂の目的である。

\end{document}
